\documentclass[border=1cm]{standalone}
\usepackage{tikz}

\begin{document}
\begin{tikzpicture}[scale=0.8]
    
    % Function to draw Sierpiński carpet
    \newcommand{\sierpinskicarpet}[4]{% x, y, size, level
        \ifnum#4=0
            \fill[black] (#1,#2) rectangle (#1+#3,#2+#3);
        \else
            \pgfmathsetmacro{\newsize}{#3/3}
            \pgfmathsetmacro{\nextlevel}{#4-1}
            
            % Draw the 8 smaller carpets around the center hole
            \sierpinskicarpet{#1}{#2}{\newsize}{\nextlevel}% bottom-left
            \sierpinskicarpet{#1+\newsize}{#2}{\newsize}{\nextlevel}% bottom-middle
            \sierpinskicarpet{#1+2*\newsize}{#2}{\newsize}{\nextlevel}% bottom-right
            
            \sierpinskicarpet{#1}{#2+\newsize}{\newsize}{\nextlevel}% middle-left
            % Center is empty - this creates the characteristic hole
            \sierpinskicarpet{#1+2*\newsize}{#2+\newsize}{\newsize}{\nextlevel}% middle-right
            
            \sierpinskicarpet{#1}{#2+2*\newsize}{\newsize}{\nextlevel}% top-left
            \sierpinskicarpet{#1+\newsize}{#2+2*\newsize}{\newsize}{\nextlevel}% top-middle
            \sierpinskicarpet{#1+2*\newsize}{#2+2*\newsize}{\newsize}{\nextlevel}% top-right
        \fi
    }

    % Draw carpets for iterations 0 to 3
    \foreach \i in {0,1,2,3} {
        \begin{scope}[xshift=\i*4cm]
            \draw[black, thick] (0,0) rectangle (3,3);
            \sierpinskicarpet{0}{0}{3}{\i}
            \node[below] at (1.5,-0.3) {Level \i};
        \end{scope}
    }
    
    % Title
    \node[above] at (6,3.5) {\Large Sierpiński Carpet};

\end{tikzpicture}
\end{document}
